%!TEX program = lualatex
\documentclass[ngerman,headheight=70pt]{scrartcl}
\usepackage[ngerman]{babel}
\usepackage{blindtext}
\usepackage[utf8]{luainputenc}
\usepackage{scrlayer-scrpage}
\usepackage[T1]{fontenc}
\usepackage{fontspec}
\usepackage{multicol}
\usepackage{enumitem}
\usepackage{epigraph}
\usepackage[german=quotes,autostyle=true,strict=true]{csquotes}
\usepackage[bottom=2cm,footskip=8mm]{geometry}
\setmainfont{Source Sans Pro}
\MakeOuterQuote{"}

\setlength{\parindent}{0pt}
\setlength{\parskip}{5pt}

\pagestyle{plain.scrheadings}
\thispagestyle{scrheadings}
\lohead{\vspace*{-1.5cm}\includegraphics[scale=0.4]{../images/stupa_logo_rot.pdf}}
\rohead{\\Präsidium des\\Studierendenparlaments}

\lofoot[Universität Hamburg · Präsidium des Studierendenparlaments\\\hrule Von-Melle-Park 5 · D-20146 Hamburg]{Universität Hamburg · Präsidium des Studierendenparlaments\\\hrule Von-Melle-Park 5 · D-20146 Hamburg}
\cofoot[]{}
\cohead[\thepage]{}

\renewcommand{\thesubsection}{TOP \arabic{subsection}}
\renewcommand{\thesubsubsection}{\arabic{subsubsection}.}
\setcounter{subsection}{-1}

\begin{document}
    UHH · StuPa-Präsidium · Von-Melle-Park 5 · D-20146 Hamburg

    \section*{Beschluss des Studierendenparlamentes vom 26. Januar 2017}
    \subsection*{Wissenschaftliche Erkenntnisse der Gesellschaft zugänglich machen,
    wissenschaftliche Arbeit besser bezahlen: Für einen VG Wort-Rahmenvertrag, der
    allen zugutekommt.}

    \subsubsection*{Was ist die VG Wort?}

    "Brotlose Kunst". Nach wie vor gelten Kunst, Kultur und Wissenschaft zwar
    als bisweilen "schön", aber doch letztendlich nutzlos, bzw. nur dann von
    Nutzen, wenn sich aus ihnen Profitables machen lässt.

    Daher liegt der Durchschnittsverdienst von Literat*innen immer noch deutlich
    unter denen von Rüstungsproduzent*innen -- im krassen Widerspruch dazu, dass
    erstere häufig Förderliches für die Menschheit produzieren, letztere nie.

    Damit dennoch unter privatwirtschaftlichen Bedingungen Kunst- und
    Kulturschaffende die materielle Grundlage ihres Lebens sicherstellen können,
    haben sich mit der dominanten Herausbildung der bürgerlichen Gesellschaft im
    18. Jahrhundert in Selbstorganisation die sogenannten "Verwertungsgesellschaften"
    konstituiert.

    Zu Ihnen gehört unter anderem die "VG Wort". Sie nimmt die Rechte der
    Autor*inen von Sprachwerken aller Art und den Verlagen wahr und trägt dafür
    Sorge, dass diese Tantiemen aus der Zweitverwertung erhalten.

    Eine solche Zweitverwertung von Sprachwerken bildet eine wesentliche Grundlage
    der Arbeit der Hochschulmitglieder insbesondere bei Lehrmaterial.

    Bislang haben die Hochschulen diese Zweitverwendung von geistiger Produktion
    pauschal mit der VG Wort abrechnen können.

    \subsubsection*{Was ist der aktuelle Konflikt?}

    Mit der Weiterentwicklung von Kopiermöglichkeiten und der Digitalisierung
    werden die Einnahmen durch die Erstverwertung für die Schaffenden der
    Schriftwerke immer geringer und die Bedeutung der Zweitverwertung steigt immer
    mehr.

    Gerade in den Hochschulen fällt dieser Trend stark ins Gewicht. Mit der
    inhaltlichen Verflachung durch das Bologna-System geht einher, dass Werke
    häufig nur noch auszugsweise in Skripten Niederschlag finden, nicht mehr aber
    in Gänze als einzelne Werke behandelt werden und daher auch keinen Eingang in
    die Bücherregale finden.

    Auch führt die Kürzungspolitik gegenüber den Hochschulen dazu, dass diese
    immer weniger in der Lage sind angemessene Pauschalen an die VG Wort zu zahlen.

    Vor diesem Hintergrund ist einzuordnen, dass die VG Wort gegen das bisherige
    Verfahren geklagt hat mit dem Erfolg eines Verfassungsgerichtsurteils, welches
    beinhaltet, dass jede einzelne Verwendung von Werken in Skripten oder in
    digitalen Semesterapparaten auch einzeln abgerechnet werden soll.

    Der Rahmenvertrag zwischen VG Wort und der Kultusministerkonferenz, der dieses
    Verfahren regeln sollte, ist nun, insbesondere weil die
    Hochschulen erheblichen Mehraufwand und dadurch entstehende Mehrkosten
    befürchteten, gescheitert.

    \subsubsection*{Wie geht es weiter? Wie soll es weitergehen?}

    Mit der Einrichtung einer neuen Verhandlungsgruppe der VG Wort und der
    Kultusministerkonferenz (KMK) besteht nun die Möglichkeit, zu einer vernünftigen
    Lösung zu kommen.

    Bis zur Klärung wird das alte Verfahren mit den Pauschalen fortgesetzt.
 
    Das Studierendenparlament hält folgende Eckpunkte einer neuen Vereinbarung
    für notwendig und fordert die Wissenschaftssenatorin aus Hamburg auf,
    entsprechend bei den Verhandlungen einzugreifen:

    \begin{itemize}
        \item Eine angemessene Vergütung, die die Lage der Autor*innen verbessert
              und nicht verschlechtert, muss sichergestellt werden.
        \item Die finanzielle Ausstattung der Hochschulen mit Grundmitteln muss
            deutlich verbessert werden, nicht nur, um die zu erwartenden Mehrausgaben
            zu decken, sondern auch, um die jahrzehntelange Unterfinanzierung zu
            beheben.
        \item Ein Verfahren sollte gefunden werden, dass für die Hochschulen
        möglichst wenig Aufwand jenseits der Kernaufgaben von Forschung, Lehre und
        Selbstverwaltung bedeutet. Diese sind bereits jetzt schon mit dem
        neoliberalen "new public management", überbordendem
        Drittmitteleinwerbungsaufwand, Evaluationshysterie und restriktiver
        Studierverhaltenskontrolle stark eingeschränkt.
        \item Da die Produktion -- auch die geistige -- prinzipiell eine
        gesellschaftliche und nie eine rein individuelle ist, ist von der VG
        Wort zu prüfen, ob bei ausreichender Gesamtbezahlung eine Pauschalabrechnung
        der Realität nicht besser gerecht wird, als eine den neoliberalen Ideologemen
        der Eigenverantwortung und Leistungsfähigkeit aufsitzende Einzelabrechnung.
    \end{itemize}

    Die Auseinandersetzung verdeutlicht über die Verhandlungen zwischen VG Wort
    und KMK hinaus folgende weitere Entwicklungserfordernisse:
    \begin{itemize}
        \item Die progressive Studienreform, die auf kritische Persönlichkeitsbildung
        in verantwortlichem Gesellschaftsbezug setzt, muss weiter vorangetrieben
        und beschleunigt werden, für exemplarisches Lernen und Vertiefung statt
        Bulimielernen und markorientierte Oberflächlichkeit. Damit wird auch die
        Würdigung von Gesamtwerken wieder gestärkt und damit die Erstverwertung
        von Schriftwerken.
        \item Das Urheberrechtsgesetz sollte dahingehend überarbeitet werden, dass
        in staatlichen und staatlich finanzierten Einrichtungen produzierte Werke
        bereits in der Produktion so ausreichend honoriert werden, dass die Erst-
        und Zweitverwertung insbesondere in öffentlichen Bildungseinrichtungen
        entfallen kann. Die in staatlich finanzierten Einrichtungen generierten
        Erkenntnisse sollten allen Mitgliedern der Gesellschaft kostenfrei
        zugänglich gemacht werden. Zugleich muss verhindert werden, dass Dritte
        mit dem kostenfreien Zugang ihrerseits wieder Geschäfte machen.
    \end{itemize}

    Sinnvolle, anregende und bereichernde wissenschaftliche und künstlerische Arbeit,
    der allgemeine Zugang zu geistigen Hervorbringungen und eine soziale Entwicklung
    zur materiellen Absicherung der Kultur- und Wissenschaftsschaffenden bilden
    eine untrennbare Einheit.
\end{document}
